% !PRINTED 569
\weberpage{569}
%<<CONT
set, wherein $\psi$ is an integral functional and $\varepsilon$ a unit, and
obtain from (14)
\weq{15}{N_a(\psi) < \sqrt{\pm 0{,}404\, \Delta}.}

Since we may set for $\varphi$ any integral functional, $\psi$ can be
a representative of an ideal of an arbitrary class, and
we have thus proved the important theorem:

3. \gesperrt{In every ideal class we can find a repre-}
\gesperrt{sentative whose norm is less than the}
\gesperrt{square root of the absolute value of the}
\gesperrt{fundamental number of the field.}

The upper limit for the norm of an ideal of a class that was
used in \S.\,153 for the proof of the finiteness of the class number
is hereby determined in a much sharper manner.

But the inequality (15) admits yet another very remarkable
conclusion. For since $N_a(\psi)$ is a positive integer, and thus
at least equal to 1, it follows from (15)
that, apart from the case $n=1$, $\Delta$ must be absolutely greater than
1, and \gesperrt{that therefore, besides the field of the rational
numbers, there is no field whose fundamental number would be $\pm 1$.} This
is the proof of this important theorem, earlier long sought
in vain, first supplied by \gesperrt{Minkowski}\footnote{In a letter to \gesperrt{Hermite} (Comptes rendus d. Pariser Akademie,
26. Januar 1891), and also in the book mentioned above, \gesperrt{Min-}
\gesperrt{kowski} finds a far narrower limit, in that he considers not only the minimum
of the quadratic forms, but that of still other functions as well,
such as $\eta_1^p + \eta_2^p + \cdots + \eta_n^p$, wherein $p$ is an arbitrary positive number, which
does not even need to be integral. For $p=1$ there results, if $n-\nu$
is the number of conjugate imaginary pairs:
\[ N_a(\psi) < \left(\frac{4}{\pi}\right)^{n-\nu} \frac{1.2.3\ldots n}{n^n}\, \sqrt{\pm \Delta}. \]
% [sic: Weber prints ``keiner'' for ``kleiner'' (less than).]
From this still further conclusions may be drawn, such as these: the fundamental
 number of a quadratic field must be greater than 4 or less than
$-2$; the fundamental number of a cubic field must be greater than 20
or less than $-12$, etc.}.

\begin{center}
\S.\,157.

\gesperrt{Prime factors of the natural primes.}
\end{center}

A more exact investigation of the basis form $\tau$ of $\mathfrak{o}$ of the
 field $\Omega$, explained in \S.\,146, shall now give us the means,
%CONT>>
